\documentclass[11pt,a4paper]{article}
\usepackage[utf8]{inputenc}
\usepackage[T1]{fontenc}
\usepackage[spanish]{babel}
\usepackage[margin=2.4cm]{geometry}
\usepackage{booktabs}
\usepackage{enumitem}
\usepackage{xcolor}
\definecolor{granate}{HTML}{7B1522}
\usepackage[colorlinks=true, urlcolor=granate, linkcolor=granate]{hyperref}
\setlength{\parskip}{4pt}

\title{\color{granate}\textbf{Metodología de búsqueda de información\\y construcción del estado del arte}}
\author{Carlos Pérez Pérez\\
\small Maestría en Inteligencia Artificial · Universidad Nacional de Ingeniería\\
\small MIA-10 Procesamiento del Lenguaje Natural · carlosperez100@gmail.com}
\date{16 de agosto de 2026}

\begin{document}
\maketitle

\noindent Este anexo documenta el método con que se identificó, seleccionó y
verificó la literatura del trabajo \emph{Detección automática de eventos
adversos hospitalarios en notas clínicas}. El artículo principal cita 34
fuentes; aquí se explica cómo se llegó a ellas.

\section{Estrategia general: cinco ejes temáticos}

La búsqueda no fue una sola consulta amplia sino \textbf{cinco búsquedas
dirigidas}, una por cada eje que la pregunta de investigación exige cubrir.
Definir primero los ejes evita el sesgo de acumular literatura de un solo
frente:

\begin{center}
\begin{tabular}{@{}p{4.6cm}p{9.6cm}@{}}
\toprule
\textbf{Eje} & \textbf{Qué debía responder la literatura} \\
\midrule
1. Epidemiología del evento adverso & ¿Cuál es la magnitud del daño y del subregistro? \\
2. PLN sobre texto clínico & ¿Qué desempeño alcanza leer notas frente a usar códigos? \\
3. Modelos de lenguaje clínicos & ¿Qué arquitecturas existen y cuál es su límite (ventana)? \\
4. Supervisión débil & ¿Cómo etiquetar corpus masivos sin anotación manual? \\
5. Validez y confiabilidad & ¿Cómo demostrar que un modelo acierta por la razón correcta y cómo medir el acuerdo humano? \\
\bottomrule
\end{tabular}
\end{center}

\section{Fuentes consultadas}

\begin{itemize}[leftmargin=1.4em, itemsep=2pt]
  \item \textbf{Literatura biomédica revisada por pares:} PubMed y Google
  Scholar (ejes 1, 2 y 5): JAMA, Health Affairs, NEJM, PLoS Medicine,
  J.~Biomedical Informatics, JAMIA.
  \item \textbf{Literatura de PLN y aprendizaje automático:} arXiv, ACL
  Anthology y actas de NAACL/EMNLP (ejes 3 y 4); Nature Machine
  Intelligence (eje 5).
  \item \textbf{Normativa institucional:} las directivas institucionales peruanas de seguridad del
  paciente (2020 y 2021, Anexos 02 y 03; referencia reservada), Plan Global de Seguridad del
  Paciente OMS 2021--2030.
  \item \textbf{Documentación técnica y repositorios:} PhysioNet (MIMIC-IV y
  su \emph{Data Use Agreement}), Hugging Face (fichas de modelos),
  scikit-learn.
  \item \textbf{Rankings especializados:} MTEB/MMTEB (pestaña de
  clasificación) e IberBench para español. Se excluyó de forma
  \emph{documentada} el Open LLM Leaderboard por estar archivado por sus
  autores.
\end{itemize}

\section{Términos de búsqueda}

Combinaciones en inglés y español, ajustadas por eje. Ejemplos
representativos:

\begin{itemize}[leftmargin=1.4em, itemsep=1pt]
  \item \texttt{"adverse event" AND "natural language processing" AND
  (detection OR surveillance)}
  \item \texttt{"discharge summaries" AND (classification OR "text mining")}
  \item \texttt{clinical BERT / ClinicalBERT / BioBERT / Longformer clinical
  / long document transformer}
  \item \texttt{"weak supervision" AND (clinical OR biomedical); Snorkel}
  \item \texttt{"shortcut learning"; "hospital-specific" confounding
  radiograph; kappa PABAK prevalence paradox}
  \item \texttt{eventos adversos subregistro; seguridad del paciente
  notificación voluntaria Perú}
\end{itemize}

\section{Bola de nieve desde artículos semilla}

Identificado un artículo central por eje, se rastrearon sus referencias
(hacia atrás) y sus citantes en Google Scholar (hacia adelante). Cadenas
principales: Brennan~1991 $\rightarrow$ de~Vries~2008 $\rightarrow$
Classen~2011 (eje 1); Chapman~2001 $\rightarrow$ Murff~2011 (eje 2);
Devlin~2019 $\rightarrow$ Lee~2020 $\rightarrow$ Alsentzer~2019
$\rightarrow$ Beltagy~2020 $\rightarrow$ Li~2022 (eje 3); Ratner~2017 (eje
4); Geirhos~2020 $\rightarrow$ Zech~2018, y Cohen~1960 $\rightarrow$
Landis--Koch~1977 $\rightarrow$ Feinstein--Cicchetti~1990 $\rightarrow$
Byrt~1993 (eje 5).

\section{Criterios de inclusión y exclusión}

\begin{center}
\begin{tabular}{@{}p{7.0cm}p{7.0cm}@{}}
\toprule
\textbf{Se incluyó} & \textbf{Se excluyó} \\
\midrule
Revisión por pares, normativa oficial o documentación técnica primaria &
Blogs, tutoriales y fuentes sin arbitraje \\
Métricas pertinentes a \emph{clasificación} de documentos &
Benchmarks de generación o razonamiento (MMLU, GPQA) \\
Estudios con estándar de referencia declarado &
Cifras sin metodología verificable \\
Rankings vigentes y mantenidos &
Rankings archivados (Open LLM Leaderboard) \\
\bottomrule
\end{tabular}
\end{center}

\section{Verificación de las citas}

Toda cifra citada en el artículo se contrastó contra su \textbf{fuente
primaria} (el resumen o el texto del estudio original), no contra citas de
segunda mano. Además, el 23 de junio de 2026 se realizó una
\textbf{auditoría de citas} registrada en la bitácora del proyecto:
verificación uno a uno de autores, año, revista y DOI de cada referencia,
con corrección de las discrepancias encontradas. Las cifras de los cuatro
estudios comparados en la Tabla~VI del artículo (Murff 2011, Classen 2011,
Zech 2018 y Li 2022) se verificaron nuevamente contra la fuente el día de la
incorporación de esa tabla.

\section{Gestión de las referencias}

Las referencias se gestionan como \texttt{bibitem} dentro del propio
documento LaTeX del artículo (formato IEEEtran), bajo control de versiones
en el repositorio público del curso:
\url{https://github.com/carlosperez100/PLN_SP}. El resultado final son
\textbf{34 referencias} distribuidas en los cinco ejes, cuya síntesis
visual es el mapa conceptual «Estado del arte: cinco caminos que desembocan
en este trabajo» incluido en la carpeta de entrega.

\end{document}
